\documentclass[a4paper,12pt,oneside]{article}
\makeatletter
\def\input@path{{../../}}
\makeatother
\usepackage{discrete-practice}

\begin{document}

\PracticeHeader{2}{2}{Связность}

\Task{Задание 1}
Найдите наименьшее и наибольшее возможное количество компонент связности
простого графа, если в нем 7 вершин и 12 ребер?

\Task{Задание 2}
Найдите наименьшее и наибольшее возможное количество вершин простого графа,
если в нем 10 ребер и 3 компоненты связности?

\Task{Задание 3}
Найдите наименьшее и наибольшее возможное количество ребер графа,
если он двудольный и в нем 23 вершины?

\Task{Задание 4}
Найдите компоненты сильной и слабой связности. Постройте конденсацию графа.
\begin{center}
  \includegraphics[width=.48\linewidth]{Practice_2_1.png}\hfill
  \includegraphics[width=.48\linewidth]{Practice_2_2.png}
\end{center}

\Task{Задание 5}
10 вершин, 5 компонент сильной связности. Оценить мин/макс количество дуг.
Граф не содержит петель и кратных дуг.

\end{document}
